% Methodology and Feature Engineering

\section{Methodology}

\subsection{Problem formulation}
Each sample of the Next Day Wildfire Spread dataset~\cite{huot2022} is a $64\times64$ patch at
1\,km resolution with 12 input channels observed on day $t$ (11 environmental channels plus the
active fire mask \texttt{PrevFireMask}) and a target fire mask \texttt{FireMask} on day $t+1$.
We treat the task as per-pixel binary classification: for every pixel, predict the probability
that it is burning on day $t+1$ using only information available on day $t$.
Pixels with label $-1$ (unknown, e.g.\ cloud cover) are excluded from training and from every metric.

\subsection{Data splits}
We use the official train/eval/test split released with the dataset
(14\,979 / 1\,877 / 1\,689 patches). Huot et al.\ assign every calendar week between 2012 and 2020
at random to train, eval or test in an 8:1:1 ratio and leave a one-day buffer between weeks, so
the splits are separated in time and a fire's day-$t$/day-$t{+}1$ pair never straddles two splits.
All features are computed inside a single patch, so no feature uses pixels from another sample. The
week-based split separates samples in time, but it does not guarantee spatial or fire-event
independence: a long-lived fire can appear in weeks assigned to different splits. We have not checked
for such overlap.
The eval split is used only for early stopping and hyperparameter selection;
the test split is scored once per final model.

\subsection{Preprocessing}
The TFRecords are parsed directly from the protobuf wire format (no TensorFlow dependency) into
arrays of shape $(N,12,64,64)$. The raw export contains physically impossible values,
so before feature construction each affected channel is clipped to fixed physical bounds
(Table~\ref{tab:clip}). Wind direction is circular, so instead of clipping it is wrapped
($\theta \bmod 360$; 0.23\% of train pixels lie outside $[0, 360)$).
The dataset does not state the patch orientation, so we verified it from the data: the arrays are
stored south-up (the row index grows northward). Steep slopes facing the last row are clearly
greener than those facing the first row (mean NDVI 0.596 vs.\ 0.555, while east- and west-facing
slopes differ by only 0.008), as expected for pole-facing slopes in the Northern Hemisphere, and only
this orientation makes next-day fire more frequent downwind than upwind in every wind-speed band
(from 1.1$\times$ in calm wind to 3.1$\times$ at 8\,m/s or more, 5\,km from the fire; Results section). All
direction-dependent features (upwind fire, wind alignment, aspect) use this orientation. The bounds are constants, not statistics fitted on the data, so clipping
introduces no leakage. No normalisation is applied for the tree models, which are invariant to
monotone rescaling of individual features.

\begin{table}[t]
\centering
\caption{Physical clipping of raw channels (train split shares of affected pixels).}
\label{tab:clip}
\begin{tabular}{lll}
\hline
Channel & Bounds & Affected pixels \\
\hline
\texttt{vs} (wind speed) & $\geq 0$ m/s & $<0.01$\% \\
\texttt{tmmn}, \texttt{tmmx} & $[253.15, 330]$ K & 0.36\% \\
\texttt{sph} & $\geq 0$ & $<0.01$\% \\
\texttt{pr} & $\geq 0$ mm & 9.78\% \\
\texttt{erc} & $\geq 0$ & 0.03\% \\
\hline
\end{tabular}
\end{table}

\subsection{Pixel tables and class imbalance}
Fire pixels are rare: on the training split only 1.1\% of labelled pixels burn on day $t+1$
(a 90:1 imbalance). Each patch is flattened into one row per pixel. For the training table we
keep every fire pixel and an independent random sample of non-fire pixels with probability chosen
to give a 5:1 non-fire to fire ratio (fixed seed), giving 3.95\,M rows.
Uniform negative subsampling shifts the predicted base rate but not the ranking of pixels,
so ranking metrics on the untouched eval and test splits remain valid.
The remaining imbalance is handled by class weighting computed from the training table only:
\texttt{scale\_pos\_weight} $= N_{\text{neg}}/N_{\text{pos}} = 5$.
Eval (7.69\,M pixels) and test (6.92\,M pixels) keep every pixel in patch order, so predictions
reshape directly into $64\times64$ probability maps.

\subsection{Models and tuning}
We train XGBoost~\cite{chen2016} and LightGBM~\cite{ke2017} on all 35 features (12 raw plus 23
engineered). For each family we run a random search of 20 configurations over the space in
Table~\ref{tab:space}. Every configuration grows up to 1\,000 trees with early stopping
(patience 50) on a fixed random sample of 1\,M labelled eval pixels, and is then scored by
PR AUC on the full eval split. The best configuration per family is evaluated once on test.

\begin{table}[t]
\centering
\caption{Random search spaces (20 configurations each).}
\label{tab:space}
\begin{tabular}{lll}
\hline
Model & Hyperparameter & Range \\
\hline
XGBoost & \texttt{max\_depth} & \{4, 6, 8, 10\} \\
 & \texttt{min\_child\_weight} & \{1, 5, 20, 100\} \\
 & \texttt{gamma} & \{0, 0.1, 1\} \\
LightGBM & \texttt{num\_leaves} & \{15, 31, 63, 127, 255\} \\
 & \texttt{min\_child\_samples} & \{20, 100, 500, 2000\} \\
 & \texttt{min\_split\_gain} & \{0, 0.1, 1\} \\
Both & \texttt{learning\_rate} & log-uniform $[0.02, 0.3]$ \\
 & \texttt{subsample} & uniform $[0.6, 1]$ \\
 & \texttt{colsample\_bytree} & uniform $[0.5, 1]$ \\
 & \texttt{reg\_lambda} & log-uniform $[0.1, 10]$ \\
\hline
\end{tabular}
\end{table}

Eval PR AUC varied little across configurations (XGBoost 0.196--0.207, LightGBM 0.189--0.203),
so the results are not sensitive to the exact hyperparameters. The selected XGBoost uses depth 6,
learning rate 0.026 and 216 trees; the selected LightGBM uses 15 leaves, at least 500 samples
per leaf, learning rate 0.058 and 294 trees. Full configurations are in
\texttt{reports/best\_params.json}.

\subsection{Evaluation}
Accuracy is uninformative at 1\% prevalence, so every model is scored with the same function on
the same pixels: PR AUC (average precision), ROC AUC, and recall at 1\% and 5\% false positive
rate, which need no decision threshold. Precision, recall and F1 are reported at the threshold
that maximises F1 on the eval split, applied unchanged to test; choosing it on test would
overstate F1. The test-optimal F1 is kept only as an upper bound.

\section{Feature Engineering}
Raw channels describe each pixel in isolation, but fire spread depends on where the fire is,
which way the wind blows relative to it, how dry the air and fuel are, and the local terrain.
We engineer 23 features in five groups (Table~\ref{tab:features}); the groups are reused for the
ablation study. Unless stated, local statistics use a $k\times k$ window with replicated edges,
\texttt{PrevFireMask} $=-1$ is treated as no fire, and wind direction $\theta$ is the direction
the wind blows \emph{from} (degrees clockwise from north).

\begin{table*}[t]
\centering
\caption{Engineered features. AUC: univariate ROC AUC on the eval split, $\max(a, 1-a)$,
sign gives direction of association. Raw channel shown for comparison.}
\label{tab:features}
\begin{tabular}{llll}
\hline
Group & Feature & Definition & AUC \\
\hline
Fire & \texttt{PrevFireMask} (raw) & fire on day $t$ & 0.593 (+) \\
 & \texttt{prev\_fire\_n3/n5/n9} & mean of fire indicator in $3/5/9$ window & 0.668 / 0.717 / 0.772 (+) \\
 & \texttt{patch\_fire\_frac} & fraction of patch burning & 0.759 (+) \\
 & \texttt{dist\_to\_fire} & Euclidean distance to nearest fire (px $\approx$ km), $64\sqrt2$ if none & 0.866 ($-$) \\
 & \texttt{upwind\_fire\_k2/k4/k8} & \texttt{prev\_fire\_n3} at the pixel $2/4/8$ px upwind & 0.649 / 0.615 / 0.586 (+) \\
 & \texttt{fire\_wind\_alignment} & $\cos$(upwind direction, direction to nearest fire) & 0.501 \\
\hline
Weather & \texttt{sph} (raw) & specific humidity & 0.565 ($-$) \\
 & \texttt{wind\_u}, \texttt{wind\_v} & $-v_s\sin\theta$, $-v_s\cos\theta$ & 0.522 / 0.518 \\
 & \texttt{temp\_range} & \texttt{tmmx} $-$ \texttt{tmmn} & 0.602 (+) \\
 & \texttt{vpd} & $e_s(T_{\max}) - e$ & 0.685 (+) \\
 & \texttt{min\_rel\_humidity} & $e / e_s(T_{\max})$ & 0.705 ($-$) \\
 & \texttt{erc\_x\_wind} & \texttt{erc} $\times v_s$ & 0.679 (+) \\
\hline
Vegetation & \texttt{NDVI} (raw) & NDVI $\times10^4$ & 0.605 ($-$) \\
 & \texttt{ndvi\_n5\_mean/std} & local mean / std of NDVI, $5\times5$ & 0.610 ($-$) / 0.513 \\
 & \texttt{ndvi\_anomaly} & NDVI $-$ local mean & 0.508 \\
\hline
Terrain & \texttt{elevation} (raw) & SRTM elevation, m & 0.626 (+) \\
 & \texttt{slope} & $\arctan\lVert\nabla z\rVert$, degrees & 0.677 (+) \\
 & \texttt{aspect\_sin/cos} & downslope direction & 0.508 \\
 & \texttt{elev\_rel} & $z - $ mean of $z$ in $9\times9$ & 0.527 (+) \\
\hline
Human & \texttt{population} (raw) & people / km$^2$ & 0.601 ($-$) \\
 & \texttt{pop\_n9} & mean population in $9\times9$ & 0.628 ($-$) \\
\hline
\end{tabular}
\end{table*}

\textbf{Fire neighbourhood.} Spread happens at the fire front, so we measure burning density
around each pixel at three scales, the distance to the nearest burning pixel, and the fraction of
the patch already burning. Three wind-aware features look for fire $k \in \{2,4,8\}$ pixels upwind (8 because 40\% of
next-day fire pixels lie more than 5\,km from day-$t$ fire, per our EDA), and an
alignment score is the cosine between the upwind direction and the direction to the nearest fire
($+1$ when the wind pushes the nearest fire toward the pixel).

\textbf{Weather.} Wind direction is circular, so it is replaced by $u$/$v$ components.
Specific humidity is converted to vapour pressure $e = qp/(0.622+0.378q)$, with surface pressure
$p$ from elevation (FAO-56). Saturation vapour pressure at the daily maximum temperature,
$e_s = 0.6108\exp\!\big(17.27T/(T+237.3)\big)$ kPa, gives the vapour pressure deficit
$e_s - e$ and minimum relative humidity $e/e_s$, both direct measures of atmospheric drying.
We also add diurnal temperature range and the interaction \texttt{erc}$\times$wind speed.

\textbf{Vegetation, terrain, human.} NDVI (stored $\times10^4$) is summarised by its local mean,
standard deviation and anomaly. Terrain adds slope and aspect from the elevation gradient and
elevation relative to the $9\times9$ neighbourhood (ridges vs.\ valleys). Population is
smoothed over $9\times9$ pixels as a proxy for settlement proximity and suppression effort.


Engineered features consistently carry more signal than the raw channels they are built from:
\texttt{dist\_to\_fire} reaches a univariate AUC of 0.866 against 0.593 for the raw fire mask,
and \texttt{min\_rel\_humidity} 0.705 against 0.565 for specific humidity.
\texttt{fire\_wind\_alignment} has no marginal signal (0.501) because it is only meaningful close
to an active fire, and gain importance in the tuned XGBoost shows the model barely uses it
(23rd of 35 features, 0.5\% of total gain). The fire-neighbourhood group dominates instead:
\texttt{dist\_to\_fire} (32\%), \texttt{prev\_fire\_n9} (23\%) and \texttt{prev\_fire\_n5} (13\%).
The three upwind features rank 4th, 5th and 7th and together hold 13.5\% of gain;
\texttt{upwind\_fire\_k8}, added after the EDA showed long-range spread, ranks 5th (3.9\%).
With the patch orientation corrected, their combined share rose from 9.4\% to 13.5\%,
so the model uses wind direction more once upwind is computed on the correct side.
The raw \texttt{PrevFireMask} ranks only 32nd because the engineered neighbourhood features already
contain its information.
